\clearpage
\addcontentsline{toc}{section}{References and audited sources}
\begin{thebibliography}{99}
\bibitem{proveit}
Vladimir Reshetnikov and ProveIt contributors.
\emph{ProveIt: Topology/UnknotRecognition}.
Active scanner, reports 06 and 07, and associated research artifacts.
Audited revision \texttt{ca61a1a5f967999d778d6efd7c882e77a2eb379d},
7 October 2026.
\url{https://github.com/VladimirReshetnikov/ProveIt/tree/ca61a1a5f967999d778d6efd7c882e77a2eb379d/Topology/UnknotRecognition}.

\bibitem{barnatan}
Dror Bar-Natan.
\emph{Fast Khovanov homology computations}.
Preprint, 2006; published in Journal of Knot Theory and Its Ramifications,
16 (2007), 243--255.
\url{https://arxiv.org/abs/math/0606318}.
Author's accompanying material:
\url{https://www.math.utoronto.ca/~drorbn/papers/FastKh/}.

\bibitem{kelomaki}
Tuomas Kelom\"aki and Dirk Sch\"utz.
\emph{On computational complexity of Khovanov homology}.
Preprint, 5 January 2026, arXiv:2601.02119v1.
In particular Theorems 1.2--1.4, Proposition 1.5, and Section 5.
\url{https://arxiv.org/abs/2601.02119v1}.

\bibitem{lackenby2026}
Marc Lackenby.
\emph{Incompressible surfaces, hierarchies and unknot recognition}.
Preprint, 25 July 2026, arXiv:2607.23350v1.
In particular Proposition 2.6, Proposition 9.1, Proposition 10.2,
Proposition 14.2, and Theorem 14.4.
\url{https://arxiv.org/abs/2607.23350v1}.

\bibitem{lackenbytalk}
Marc Lackenby.
\emph{Unknot recognition in quasi-polynomial time}.
Public slides, February 2021, hosted by the
University of Oxford.  Slides 13--14 (PDF pages counted from one)
display $n^{O(\log n)}$.
\url{https://people.maths.ox.ac.uk/lackenby/quasipolynomial-talk.pdf}.

\bibitem{lackenby2024}
Marc Lackenby.
\emph{Unknot recognition is quasi-polynomial time}.
Seminar announcement and abstract, Yau Mathematical Sciences Center,
Tsinghua University, 7 November 2024 (announcement dated 6 November).
The abstract displays $2^{O((\log n)^3)}$.
\url{https://qzc.tsinghua.edu.cn/en/info/1122/4171.htm}.

\bibitem{openai_artin}
OpenAI.
\emph{Parabolic Intersections in Artin Groups}.
Repository manuscript, 23 September 2026, \texttt{openai/math};
audited revision \texttt{adc7f1241b42e322a6451854ab7e4b4c146bf78a}.
The source reviewed here is the lemma ``Composition at the bottom'' and
its bounded-layer construction, together with the categorical-action
section; no other claim of the manuscript is required.
\url{https://github.com/openai/math/blob/adc7f1241b42e322a6451854ab7e4b4c146bf78a/preprints/Parabolic-intersections-in-Artin-groups-September-23-2026/build/sections/03_layers.tex}.
Main source: \texttt{build/paper.tex} in the same manuscript directory.

\bibitem{kronheimer}
P. B. Kronheimer and T. S. Mrowka.
\emph{Khovanov homology is an unknot-detector}.
Publications Math\'ematiques de l'IH\'ES 113 (2011), 97--208;
preprint arXiv:1005.4346, 2010.
\url{https://arxiv.org/abs/1005.4346}.

\bibitem{shumakovitch}
Alexander N. Shumakovitch.
\emph{Torsion of the Khovanov homology}.
Preprint arXiv:math/0405474, 2004; version 2, 19 June 2018.
Corollary 3.2.C supplies the reduced/unreduced splitting over
$\mathbb F_2$.
\url{https://arxiv.org/abs/math/0405474v2}.

\bibitem{crainic}
Marius Crainic.
\emph{On the perturbation lemma, and deformations}.
Preprint, 2004, arXiv:math/0403266.
\url{https://arxiv.org/abs/math/0403266}.

\bibitem{khovanov_arc}
Mikhail Khovanov.
\emph{A functor-valued invariant of tangles}.
Preprint, 2001, arXiv:math/0103190.
\url{https://arxiv.org/abs/math/0103190}.
\end{thebibliography}
